% Seitenrand: Innen mind. 25mm
% Außenrand: Korrekturrand 50mm,
%   Randnotizen können in den 50mm-Rand hineinragen

Die Seitenränder sollten gemäß DIN~5008 mind. \qty{2.5}{\centi\metre} betragen.
Bei Arbeiten, die auf der Außenseite einen Korrekturrand erfordern, sollte der
Korrekturrand \qty{5}{\centi\metre} breit sein.

Die Festlegungen zur Seitengeometrie werden in der Drucksprache als
"`Satzspiegel"' bezeichnet. Im Laufe der Geschichte des Buchdrucks haben sich
viele Regeln zur Gestaltung eines guten Satzspiegels eingebürgert. Wenn Sie
einmal nach den Stichwörtern "`Satzspiegel"' und "`\LaTeX"' recherchieren,
werden Sie viele Abhandlungen finden. Halten Sie sich im Zweifel aber eher
an die Vorgaben des Dozenten, auch wenn diese den klassischen Regeln der
Satzspiegelkonstruktion widersprechen.

Die Einstellung des Satzspiegels erfolgt mit dem Paket~\ctanref{geometry}.
Listing~\ref{lst:tpl:margin} enthält ein Beispiel für die Anwendung des Pakets.
Die Option \texttt{includeheadfoot} besagt, dass der obere bzw. untere Rand bis
zur Oberkante der Kopfzeile bzw. bis zur Unterkante der Fußzeile gemessen wird.
Andernfalls wird der Rand bis zum Textkörper gemessen.

\lstinputlisting[language={[LaTeX]{TeX}},style=block,float,caption={Einstellung der Seitenränder},label={lst:tpl:margin}]{code/tpl_margin.de.tex}

Das Paket~\texttt{geometry} bietet noch weitaus mehr Möglichkeiten, die den
Rahmen dieser Einführung sprengen würden. Hierfür sei auf die
Paketdokumentation verwiesen.


\begingroup
\locbox\evenlayout
\locbox\oddlayout

\printheadingsfalse
\printparametersfalse
%\drawdimensionstrue
\currentpage
\setlayoutscale{0.45}

\oddpagelayoutfalse\savebox{\evenlayout}{{\hsize=0.46\paperwidth\vbox{\centering\footnotesize\pagedesign}}}
\oddpagelayouttrue\savebox{\oddlayout}{{\hsize=0.46\paperwidth\vbox{\centering\footnotesize\pagedesign}}}

\begin{sidewaysfigure*}[p]
\begin{minipage}{0.5\linewidth}
\centering
\usebox{\evenlayout}
\end{minipage}%
\begin{minipage}{0.5\linewidth}
\centering
\usebox{\oddlayout}
\end{minipage}
\caption{Satzspiegel dieses Dokuments}
\label{fig:tpl:margin_layout}
\end{sidewaysfigure*}
\endgroup

Durch Einfügen des Pakets~\ctanref{showframe} können Sie sich den Satzspiegel
Ihres Dokuments in einer Form vergleichbar zu
Abbildung~\ref{fig:tpl:margin_layout} anzeigen lassen.

%\lstinputlisting[language={[LaTeX]{TeX}},style=block,float,caption={Anzeige des Satzspiegels},label={lst:tpl:margin_showframe}]{code/tpl_margin_showframe.de.tex}
